% Recovery reconstruction after workspace loss.  The original native body had
% SHA256 21bea3e64e4ae625ed00dd5a94027d2738e921e7cd6f8016b8f8f2c279f79b48.
% This file preserves its proved statements and constants; byte identity is
% not asserted.  Fresh compilation and mathematical checks are required.
\section{Boundary localization of a global disk-word minimum}
\label{sec:bl-global}
The summable sublevel theorem gives a collision-independent budget,
but does not by itself locate its mass.  Here global minimality pays
the variational work in the interior by work in two prescribed
transition layers.  A nonlinear recurrence then gives a uniform
algebraic tail.  No regular physical continuation from the reference
to the comparison is assumed.

Let $m\ge1$ and let $Q_0,\ldots,Q_m$ be an original regular physical record, or a
fixed-word closed physical reference, at an original radius
$R\in[R_-,r]=[9/20,47/100]$.  Its ordered centers are $O_i$.
Let $Y$ be any global minimum of the unchanged disk-word action
\[
 A(Y)=\sum_{i=0}^{m-1}|Y_{i+1}-Y_i|,
 \qquad Y_i\in\overline B(O_i,R).
\]
Consecutive centers are distinct.  All original and comparison
flights therefore have length at least $\tau_-=3/50$.  Use the
retained physical constants
\[
 N=94,\qquad a=1/416640,\qquad T=187/25,\qquad C=10^8.
\]
Every complete $N$-flight reference block has an interior contact
with incidence at least $a$.  The retained physical horizon is
$\tau_i\le28/5$; hence every feasible comparison edge satisfies
$\ell_i\le\tau_i+|\delta_i|+|\delta_{i+1}|\le28/5+4r=T$.
These statements include fixed-word closed limits.
Write
\begin{equation}\label{eq:bl-notation}
 \begin{gathered}
 \delta_i=Y_i-Q_i,\quad
 \tau_i=|Q_{i+1}-Q_i|,\quad \ell_i=|Y_{i+1}-Y_i|,\\
 v_i=(Q_{i+1}-Q_i)/\tau_i,\quad
 \widehat v_i=(Y_{i+1}-Y_i)/\ell_i,\quad e_i=\widehat v_i-v_i,\\
 F_i=\ell_i|e_i|^2,\qquad
 H_i=c_i|\delta_i|^2\quad(0<i<m),\qquad
 E=E_R(Q)=R(2-c_0-c_m).
 \end{gathered}
\end{equation}
The reference dual certificate $D_Q$ gives the complete positive
gap $G=A(Y)-D_Q$.  The inherited gap identity and disk support
inequality give
\[
 0\le G\le E\le2r,\qquad
 G=\tfrac12\sum_iF_i+\sum_{i=0}^m h_i,
 \qquad h_i\ge H_i/R\ (0<i<m).
\]
The already proved sublevel theorem gives
$\sum_{i=0}^m|\delta_i|^4\le CE$.  We do not reclassify that
theorem as a contribution of this section.

\begin{lemma}[A variational inequality with only cutoff work]
\label{lem:bl-caccioppoli}
For any $0\le\eta_i\le1$ with $\eta_0=\eta_m=0$, put
$a_i^\eta=(\eta_i^2+\eta_{i+1}^2)/2$.  Then
\begin{equation}\label{eq:bl-caccioppoli}
 \frac1R\sum_{i=1}^{m-1}\eta_i^2H_i
 +\frac14\sum_{i=0}^{m-1}a_i^\eta(\ell_i+\tau_i)|e_i|^2
 \le\frac1{\tau_-}\sum_{i=0}^{m-1}
 (\eta_{i+1}-\eta_i)^2(|\delta_i|^2+|\delta_{i+1}|^2).
\end{equation}
In particular, no separate endpoint deficit is paid on the right.
\end{lemma}
\begin{proof}
Let $n_i=(Q_i-O_i)/R$.  At an interior reference contact the
action gradient is $g_i^Q=v_{i-1}-v_i=-2c_in_i$.
For every point in the same disk,
\[
 |\delta_i|^2\le-2R n_i\cdot\delta_i,
 \qquad g_i^Q\cdot\delta_i\ge H_i/R.
\]
The one-coordinate minimum condition gives
$g_i^Y\cdot(Q_i-Y_i)\ge0$.  Since
$g_i^Q-g_i^Y=e_i-e_{i-1}$, it follows that
$(e_i-e_{i-1})\cdot\delta_i\ge H_i/R$.
This argument applies to inactive and interior comparison vertices.

Set $b_i^\eta=\eta_{i+1}^2-\eta_i^2$ and
$\overline\delta_i=(\delta_i+\delta_{i+1})/2$.  The exact identities
\[
 \eta_{i+1}^2\delta_{i+1}-\eta_i^2\delta_i
 =a_i^\eta(\delta_{i+1}-\delta_i)+b_i^\eta\overline\delta_i,
 \qquad
 e_i\cdot(\delta_{i+1}-\delta_i)
       =\tfrac12(\ell_i+\tau_i)|e_i|^2
\]
and finite integration by parts put the contact work and the full
kinetic term on the left.  Their sum is at most
$-\sum_i b_i^\eta e_i\cdot\overline\delta_i$.
When $a_i^\eta>0$, Young's inequality bounds the absolute value
of each summand by
\[
 \tfrac14a_i^\eta(\ell_i+\tau_i)|e_i|^2+
 \frac{(b_i^\eta)^2|\overline\delta_i|^2}
                         {a_i^\eta(\ell_i+\tau_i)}.
\]
For $a_i^\eta=0$ both associated cutoff values and the work vanish.
Use
\[
 (b_i^\eta)^2/a_i^\eta\le4(\eta_{i+1}-\eta_i)^2,
 \quad |\overline\delta_i|^2
       \le(|\delta_i|^2+|\delta_{i+1}|^2)/2,
 \quad \ell_i+\tau_i\ge2\tau_-.
\]
Absorbing the first Young term proves the stated coefficient and
normalization.  Only the global minimum condition was used for
$Y$; the inequality is not asserted for an arbitrary comparison.
\end{proof}

\begin{lemma}[Interior coercivity without endpoint payment]
\label{lem:bl-interior}
For $0\le u<v\le m$, define
\begin{equation}\label{eq:bl-local-constants}
 A_H=\frac{4r^2}{a}
       +8r^2\left(1+\frac{3N}{a}+12N^2\right),
 \qquad A_F=24r^2N^2T.
\end{equation}
Then
\begin{equation}\label{eq:bl-local-coercivity}
 \sum_{i=u+N}^{v-1}|\delta_i|^4
 \le A_H\sum_{u<i<v}H_i+A_F\sum_{u\le i<v}F_i.
\end{equation}
An empty sum is zero.  Every anchor used to prove this estimate
is strictly inside $[u,v]$.  This lemma holds for every feasible
comparison polygon, without a minimum assumption.
\end{lemma}
\begin{proof}
For an included contact write
$b_i=v_{i-1}^{\perp}\cdot\delta_i$.  At an anchor,
$|\delta_i|^4\le4r^2H_i/a$.
At any other included contact choose the last preceding anchor
$j$; the complete $N$-flight block preceding $i$ supplies one
with $u<j<i$ and $i-j<N$.  The exact transverse telescope is
\[
 b_i=b_j+\sum_{l=j}^{i-1}\ell_l v_l^\perp\cdot e_l
          +\sum_{l=j}^{i-1}(v_l-v_{l-1})^\perp\cdot\delta_l.
\]
Here $|v_l-v_{l-1}|=2c_l$ and $c_l^2\le c_l$.  Cauchy--Schwarz
and the three-square inequality consequently give
\[
 b_i^2\le\frac{3H_j}{a}
             +3NT\sum_{l=j}^{i-1}F_l
             +12N\sum_{l=j}^{i-1}H_l.
\]
The same disk support inequality used above, resolved into the
incoming direction and its perpendicular, implies
\[
 |\delta_i|^2/(2R)\le c_i|\delta_i|+|b_i|,
 \qquad |\delta_i|^4\le8r^2(H_i+b_i^2).
\]
Every preceding anchor, flight, and contact occurs in at most $N$
of the chosen intervals.  Sum the anchor bounds separately, and
sum the last two displays over the nonanchors.  The resulting
coefficients are exactly $A_H,A_F$.  The indices range from
$u+N$ to $v-1$, so no endpoint surrogate enters the telescope.
\end{proof}

Put $S_d=\sum_{i=d}^{m-d}|\delta_i|^4$ for integer $d\ge0$,
with $S_d=0$ if $2d>m$.
\begin{theorem}[A nonlinear depth recurrence for every global minimum]
\label{thm:bl-recurrence}
For integers $k\ge0$ and $L\ge1$,
\begin{equation}\label{eq:bl-recurrence}
 S_{k+L+N}\le C_{\rm cut} L^{-3/2}\sqrt{S_k},
 \qquad C_{\rm cut}=\frac{4(rA_H+4A_F)}{\tau_-}.
\end{equation}
The same constant works for every original collision count, word,
radius, and global minimum.
\end{theorem}
\begin{proof}
There is nothing to prove when the target is empty.  Otherwise choose
\[
 \eta_i=\min\{1,(i-k)_+/L,(m-k-i)_+/L\}.
\]
There are two complete transition layers of $L$ flights.  At most
$4L$ vertices occur in them, each at most twice, and all have
depth at least $k$.  Therefore
\[
 \sum_{\eta_{i+1}\ne\eta_i}
       (|\delta_i|^2+|\delta_{i+1}|^2)
 \le4\sqrt{L S_k}.
\]
The right side of \eqref{eq:bl-caccioppoli} is thus at most
$W=4\tau_-^{-1}L^{-3/2}\sqrt{S_k}$.  On
$[u,v]=[k+L,m-k-L]$ the cutoff is one.  Its contact sum is at
most $rW$, and its flight sum is at most $4W$, since the retained
kinetic term is at least $F_i/4$ there.  The target positions lie
in $[u+N,v-1]$.  Lemma~\ref{lem:bl-interior} proves the claimed
recurrence with exactly the displayed $C_{\rm cut}$.
\end{proof}

\begin{theorem}[Uniform algebraic localization through inactive contacts]
\label{thm:bl-tail}
With
\begin{equation}\label{eq:bl-K-tail}
 K_{\rm tail}=\max\{2rC(4N)^3,(64C_{\rm cut})^2\},
\end{equation}
every global minimum satisfies
\begin{equation}\label{eq:bl-tail}
 S_d\le\min\{CE_R(Q),K_{\rm tail}/(d+1)^3\}
                    \qquad(d\ge0).
\end{equation}
The estimate includes an isolated central contact when $m=2d$.
\end{theorem}
\begin{proof}
The first bound is the inherited full sublevel estimate.  For
$d<4N$ it also proves the second bound by the first term in
\eqref{eq:bl-K-tail}.  For $d\ge4N$ use strong induction, put
$k=\lfloor d/2\rfloor$ and $L=d-k-N$, and note that
\[
 L\ge1,\qquad L(k+1)\ge(d+1)^2/16.
\]
The recurrence and the induction hypothesis give
\[
 S_d\le\frac{C_{\rm cut}\sqrt{K_{\rm tail}}}
                       {[L(k+1)]^{3/2}}
 \le\frac{64C_{\rm cut}\sqrt{K_{\rm tail}}}{(d+1)^3}
 \le\frac{K_{\rm tail}}{(d+1)^3}.
\]
If the target is empty the same inequality is immediate.  All
constants were fixed before either the word or the depth was chosen.
\end{proof}

\begin{corollary}[Deep contact and flight energy]
\label{cor:bl-energy}
For every integer $d\ge1$,
\begin{align}
 \sum_{\substack{0<i<m\\d\le i\le m-d}}H_i
 &\le\frac{32r\sqrt{K_{\rm tail}}}{\tau_-d^3},
                                      \label{eq:bl-contact-tail}\\
 \sum_{i=d}^{m-d-1}F_i
 &\le\frac{128\sqrt{K_{\rm tail}}}{\tau_-d^3}.
                                      \label{eq:bl-flight-tail}
\end{align}
\end{corollary}
\begin{proof}
Use the same cutoff with $k=\lfloor d/2\rfloor$ and $L=d-k$,
without an anchor offset.  Since $L(k+1)\ge d^2/4$, its work is
at most $32\sqrt{K_{\rm tail}}/(\tau_-d^3)$.
The contact and flight bounds follow respectively by multiplying
this work by $r$ and by $4$.  For a central singleton the contact
is still in the unit-cutoff region, and the flight sum is empty.
\end{proof}
